\section{Chronological products under concatenation}
\label{sec:contraction_word_append}

\begin{theorem}[Later factors act on the left]
\label{thm:contraction_word_append}
\lean{Matrix.contractionWord_append, Matrix.contractionWord_word_append}
\leanok
\uses{def:contraction_word, def:poisson_word_append}
For the chronological square-root products of
Definition~\ref{def:contraction_word}, concatenation satisfies
\begin{align}
 D_{uv}&=D_vD_u.
 \label{eq:contraction_word_append}
\end{align}
The identity depends only on the ordered matrix product, without a
commutation, order, ground-state, or gap hypothesis on the matrices $k_i$.
It holds for tuples of fixed lengths and for arbitrary finite words,
including either empty word.
\end{theorem}
\begin{proof}
\leanok
\uses{def:contraction_word}
Induct on the length of $v$, appending its last letter $i$.
The empty word gives $D_{u\varnothing}=D_u=ID_u$.
For the induction step,
$D_{u(vi)}=G_iD_{uv}=G_i(D_vD_u)=(G_iD_v)D_u=D_{vi}D_u$.
Thus the factors of the later word multiply on the left, as in
\cite{OpenAI2026AreaLaw}, \texttt{09-amplification.tex}, lines 49--54.
\end{proof}

Together with~(\ref{eq:poisson_word_append_law}), this identifies the
matrix-product orientation for distributional time composition.
It supplies no event times, process with independent increments,
path coupling, or locality assertion.
